\documentclass[10pt,a4paper]{amsart}
\usepackage[margin=19mm]{geometry}
\usepackage{amsmath,amssymb,mathtools}
\usepackage[scr=rsfs,cal=euler]{mathalfa}
\usepackage{xfrac}
\usepackage[unicode,hidelinks,bookmarks=false]{hyperref}
\newcommand{\ie}{i.e.\ }
\newcommand{\cf}{cf.\ }
\newcommand{\resp}{respectively\ }
\newcommand{\Subsection}{\subsection}
\providecommand{\nobreakdash}{\nobreak\hbox{-}}
\setlength{\emergencystretch}{2em}
\multlinegap0pt
\usepackage{amssymb}
\usepackage{csquotes}
\usepackage[shortlabels]{enumitem}
\usepackage{tikz}
\newtheorem{prop}{Proposition}
\title{The weak Banach-Saks property for Hölder spaces}
\author{Prezemys{\l}aw G\'orka, Mauro Sanchiz}
\date{Source: arXiv:2603.02977v1 (CC BY 4.0)}
\begin{document}
\maketitle

\noindent\emph{Source and licence.} The weak Banach-Saks property for Hölder spaces. Version arXiv:2603.02977v1 (CC BY 4.0), distributed by arXiv under the Creative Commons Attribution 4.0 International licence, http://creativecommons.org/licenses/by/4.0/, as recorded in the arXiv licence metadata for this exact version and shown on the abstract page https://arxiv.org/abs/2603.02977v1. Full licence notice: \texttt{licenses/arxiv-2603.02977-CC-BY-4.0.txt}.\par\smallskip

\noindent\textit{Editorial scope.} Counted unit: the article's prop iso (source0.tex lines 249--251). The article's complete original proof of it is delivered with it (source0.tex lines 252--304, one \texttt{proof} environment of 53 lines). No mathematics is written, repaired or re-derived: the statement and the proof are the article's own lines, in the article's own notation, and no retained line is substituted or re-flowed.  No further source-local context is needed: every cross-reference of every retained line resolves inside the two delivered spans.  Nothing was omitted from any retained span, so the package carries no omission record.  No retained line contains a citation, so the wrapper carries no bibliography and no bibliography entry.  No retained line contains a figure environment, an image-inclusion command or drawing code, so no image is delivered and the package declares no asset dependency.  Editorial formatting only, in addition: an AMS \texttt{amsart} preamble that restates the article's own declarations and macros used by the retained lines, namely the environments \texttt{prop} on separate counters, together with the standard packages those lines need.  The retained environments print in this document as Proposition 1 in order of appearance, whereas the article prints the same items with the numbering of its own full document.  The complete native source of the article as distributed in the arXiv e-print for this exact version is bundled unchanged as \texttt{sources/195646/source0.tex}.
\begin{prop} \label{iso}
    Let $(M,d)$ be a metric space, and let $(\widehat{M},\widehat{d})$ be its completion. Fix $\alpha \in (0,1]$. Then the following $C^\alpha(\widehat{M}) \cong C^\alpha(M)$ isometric isomorphism holds.
\end{prop}

\begin{proof}
Let $(\widehat{M},\widehat{d})$ denote the completion of $(M,d)$, and let 
$i: M \to \widehat{M}$ be an isometric embedding with dense image. 
Define the operator
\[
I: C^{\alpha}(\widehat{M}) \to C^{\alpha}(M)
\]
by
\[
I(f)(x) = f(i(x)), \qquad \text{for all } f \in C^{\alpha}(\widehat{M}),\ x \in M.
\]

It is immediate that $I$ is linear and injective. Moreover, for every 
$f \in C^{\alpha}(\widehat{M})$,
\[
\|I(f)\|_{C^{\alpha}(M)} \le \|f\|_{C^{\alpha}(\widehat{M})}.
\]

We now prove that $I$ is surjective. Let $g \in C^{\alpha}(M)$ be given and define 
$\widetilde{f}: i(M) \to \mathbb{R}$ by
\[
\widetilde{f}(\hat{x}) = g(i^{-1}(\hat{x})).
\]
For any $\hat{x}, \hat{y} \in i(M)$ we obtain
\[
|\widetilde{f}(\hat{x}) - \widetilde{f}(\hat{y})|
= |g(i^{-1}(\hat{x})) - g(i^{-1}(\hat{y}))|
\le \rho_{\alpha,M}(g)\,d^{\alpha}(i^{-1}(\hat{x}),i^{-1}(\hat{y}))
=\rho_{\alpha,M}(g)\,\widehat{d} ^{\alpha}(\hat{x},\hat{y}),
\]
and
\[
|\widetilde{f}(\hat{x})| \le \|g\|_{C(M)}.
\]
Thus, $\widetilde{f}$ is Hölder continuous on the dense subset 
$i(M) \subset \widehat{M}$ with the same Hölder seminorm and supremum bound as $g$.
Since $i(M)$ is dense in $\widehat{M}$, the function $\widetilde{f}$ admits 
a unique continuous extension $f: \widehat{M} \to \mathbb{R}$. 
Furthermore, this extension satisfies
\[
\rho_{\alpha,\widehat{M}}(f) \le \rho_{\alpha,M}(g),
\qquad
\|f\|_{C(\widehat{M})} \le \|g\|_{C(M)}.
\]
Hence $f \in C^{\alpha}(\widehat{M})$ and, by construction, $I(f)=g$. 
Consequently,
\[
\|f\|_{C^{\alpha}(\widehat{M})}
\le
\|I(f)\|_{C^{\alpha}(M)}.
\]
This completes the proof.
\end{proof}

\end{document}
